\documentclass[11pt,a4paper]{article}
\usepackage[utf8]{inputenc}
\usepackage[spanish,es-nodecimaldot]{babel}
\usepackage[margin=2.2cm]{geometry}
\usepackage{helvet}
\renewcommand{\familydefault}{\sfdefault}
\usepackage{amsmath,amssymb}
\usepackage{graphicx}
\usepackage{booktabs}
\usepackage{xcolor}
\usepackage{hyperref}
\usepackage{fancyhdr}
\usepackage{enumitem}
\usepackage{tcolorbox}
\usepackage{colortbl}

% Color Definitions
\definecolor{navyblue}{RGB}{20, 45, 85}
\definecolor{alertred}{RGB}{215, 38, 56}
\definecolor{alertwarning}{RGB}{245, 158, 11}
\definecolor{alertgreen}{RGB}{16, 185, 129}
\definecolor{lightbg}{RGB}{245, 247, 250}
\definecolor{accentblue}{RGB}{41, 128, 185}

\hypersetup{
    colorlinks=true,
    linkcolor=navyblue,
    urlcolor=accentblue,
    pdftitle={Reporte de Evaluación Técnica - Proyecto Júpiter},
    pdfauthor={Proyecto Júpiter - Inteligencia Hidrometeorológica}
}

% Header & Footer
\setlength{\headheight}{14.5pt}
\pagestyle{fancy}
\fancyhf{}
\rhead{\textcolor{navyblue}{\textbf{Proyecto Júpiter — Informe Técnico de Evaluación}}}
\lhead{\textcolor{gray}{La Serena / Región de Coquimbo}}
\rfoot{\thepage}
\lfoot{\textcolor{gray}{Documento de Evaluación para Bomberos y Municipalidad}}

\begin{document}

% Title Block
\begin{tcolorbox}[colback=navyblue,colframe=navyblue,arc=4mm,boxrule=0pt,top=6mm,bottom=6mm]
\centering
\color{white}
{\Huge \textbf{PROYECTO JÚPITER}}\\[2mm]
{\Large \textbf{Sistema de Alerta Temprana e Inteligencia Hidrometeorológica}}\\[3mm]
{\small \textbf{Informe de Evaluación Técnica Post-Prueba Real, Lecciones Aprendidas y Escalabilidad}}\\[1mm]
{\footnotesize Dirigido a: Equipo de Desarrollo, Cuerpo de Bomberos de La Serena e Ilustre Municipalidad de La Serena}
\end{tcolorbox}

\vspace{4mm}

\begin{tcolorbox}[colback=lightbg,colframe=accentblue,arc=2mm,title=\textbf{Resumen Ejecutivo}]
El presente documento constituye la evaluación exhaustiva de la \textbf{Fase de Prueba Real en Terreno} del Proyecto Júpiter (realizada entre el 30 de Julio y el 01 de Agosto de 2026 en La Serena). Se analizan los puntos ciegos iniciales detectados durante el evento, las refactorizaciones matemáticas y tecnológicas implementadas en tiempo real, los resultados de precisión alcanzados y la propuesta de escalabilidad institucional para su presentación ante el Cuerpo de Bomberos y la Municipalidad de La Serena.
\end{tcolorbox}

\section{Introducción y Contexto Operativo}
La comuna de La Serena presenta una compleja geografía hidro-meteorológica caracterizada por cuencas orográficas de rápida respuesta en precordillera ($T_c \approx 1.5\text{h}$) y colectores bajos urbanos costeros ($T_c \approx 4.0\text{h}$). Ante eventos de precipitación convectiva o precipitaciones orográficas en ladera, los sistemas de alerta comunales tradicionales suelen generar dos problemas críticos:
\begin{enumerate}[leftmargin=*]
    \item \textbf{Falsas Alarmas Comunales:} Declarar Alerta Roja para toda la ciudad cuando la lluvia intensa solo afecta a quebradas precordilleranas.
    \item \textbf{Puntos Ciegos de Escorrentía:} Omitir la detección de bajadas de agua e inundaciones en sectores como Pueblo Islón o Las Compañías cuando la estación meteorológica costera registra llovizna débil.
\end{enumerate}

El \textbf{Proyecto Júpiter} se diseñó para resolver esta problemática mediante Inteligencia Artificial, Asimilación de Datos EnKF, Dispersión Espacial IDW e Ingesta NLP de Noticias.

\section{Análisis de Puntos Ciegos Iniciales y Fallas Detectadas}
Durante las primeras horas de la fase de prueba real, se identificaron seis limitaciones críticas que fueron corregidas sistemáticamente:

\begin{table}[h!]
\centering
\small
\begin{tabular}{p{3.5cm} p{5.5cm} p{6.5cm}}
\toprule
\textbf{Punto Crítico Detectado} & \textbf{Síntoma Observado en Terreno} & \textbf{Causa Raíz Identificada} \\
\midrule
\rowcolor{lightbg}
\textbf{1. Falsa Alerta Roja por EnKF No Acotado} & Salto abrupto de 8.7\% a 88.1\% en el Centro Histórico con solo 2.5 mm de lluvia. & La variable de innovación $EnKF$ se sumaba sin normalizar a una escala $[0,1]$. \\
\textbf{2. Punto Ciego Fluvial en Pueblo Islón} & Crecida de río y desborde registrado sin anticipación en el modelo (mostraba 25\%). & La escorrentía solo medía lluvia local ($direct\_Q$), ignorando el caudal fluvial de aguas arriba. \\
\rowcolor{lightbg}
\textbf{3. Vulnerabilidad ante Caída de CEAZAMET} & El servidor de estaciones CEAZAMET quedó fuera de línea durante el pico de la tormenta. & Dependencia de una sola fuente primaria de terreno sin factor de incertidumbre adaptativo. \\
\textbf{4. Lluvia Plana Comunal} & Imposibilidad de diferenciar la lluvia orográfica de la ladera con la llovizna costera. & Asignación de un único escalar de precipitación comunal a los 35 micro-sectores. \\
\rowcolor{lightbg}
\textbf{5. Desconexión con Anegamiento Urbano} & Calles anegadas en Las Compañías y Ruta 5 mientras sensores marcaban baja lluvia. & Falta de ingesta de reportes ciudadanos y noticias de emergencias en tiempo real. \\
\textbf{6. Falsa Proyección de Picos Pasados} & Proyección de picos de lluvia futuros cuando el pico de la tormenta ya había transcurrido. & El algoritmo solo escaneaba datos futuros ($t \ge t_{actual}$), ignorando el hidrograma pasado. \\
\bottomrule
\end{tabular}
\caption{Matriz de Fallas Iniciales y Diagnóstico Técnico durante la Prueba Real.}
\end{table}

\section{Refactorizaciones Matemáticas e Innovaciones Implementadas}
Para resolver definitivamente las fallas detectadas, se implementó una arquitectura de 6 capas de protección:

\subsection{1. Ponderación Convexa y Suavizado Adaptativo EMA}
Se estructuró una formulación convexa donde la suma de pesos es estrictamente igual a 1.00 ($\sum w_i = 1.00$):
\begin{align}
R_{raw} &= 0.22 \cdot \phi_{precip} + 0.18 \cdot \phi_{API} + 0.15 \cdot \phi_{ML} + 0.15 \cdot \phi_{runoff} + 0.10 \cdot \phi_{forecast} \nonumber \\
        &\quad + 0.05 \cdot \phi_{geo} + 0.05 \cdot \phi_{rout} + 0.05 \cdot \phi_{EnKF} + 0.05 \cdot \phi_{news}
\end{align}
Para evitar oscilaciones o \textit{flickering}, se integró una Media Móvil Exponencial Asimétrica (EMA) con histeresis del 8\%:
\begin{equation}
S_{t} = \alpha \cdot S_{inst} + (1 - \alpha) \cdot S_{t-1}, \quad \text{donde } \alpha = \begin{cases} 0.45 & \text{si } S_{inst} \ge S_{t-1} \text{ (Adaptación rápida a subidas)} \\ 0.10 & \text{si } S_{inst} < S_{t-1} \text{ (Decaimiento paulatino en drenaje)} \end{cases}
\end{equation}

\subsection{2. Dispersión Espacial IDW (Haversine Inverse Distance Weighting)}
Se calibró un motor de interpolación espacial por micro-sectores que calcula la precipitación $P_{sector}$ de cada una de las 35 zonas según las coordenadas WGS84 de las estaciones CEAZAMET y la grilla satelital:
\begin{equation}
P_{sector}(lat_i, lon_i) = \frac{\sum_{j=1}^{K} w_j \cdot P_j}{\sum_{j=1}^{K} w_j}, \quad w_j = \frac{1}{d(lat_i, lon_i, lat_j, lon_j)^2}
\end{equation}

\subsection{3. Ingesta Redundante de 5 Proveedores Meteorológicos}
Se desarrolló un módulo ingesdor redundante que consulta en paralelo:
\begin{itemize}[leftmargin=*]
    \item Open-Meteo Current Telemetry API.
    \item Open-Meteo Multi-Model Ensemble (ECMWF IFS, GFS Seamless, ICON Seamless).
    \item wttr.in International Weather Service.
    \item OpenWeatherMap GFS Telemetry.
    \item MeteoBlue / Tomorrow.io Contingency Fallback.
\end{itemize}
Si una fuente cae, el sistema aplica un factor de incremento por incertidumbre ($Uncertainty\_Boost = +15\%$).

\subsection{4. Ingestor NLP de Noticias de Terreno y Alertas SAE}
Se creó el módulo \texttt{ingest\_news\_reports.py} que procesa titulares de medios regionales (El Día, Diario La Serena, BioBioChile) y alertas oficiales de SENAPRED. Si se detectan palabras clave de anegamiento en Las Compañías o Ruta 5, el modelo eleva el score base a mínimo 70.0\% (\textbf{Alerta Roja Focalizada}).

\subsection{5. Motor de Detección de Picos Transcurridos e Hidrograma}
El sistema analiza el hidrograma en una ventana de 24 horas ($t_{-12h}$ a $t_{+12h}$). Si el tiempo del pico máximo ($t_{pico} + T_c$) es menor a la hora actual, el modelo clasifica automáticamente el estado como \textbf{"punto máximo transcurrido"} y despliega la cuenta regresiva para el \textbf{Paso Seguro por Drenaje Expirado}.

\section{Resultados Alcanzados en Tiempo Real}
Durante la prueba real, el modelo demostró una alta precisión táctica y coherencia física:

\begin{table}[h!]
\centering
\small
\begin{tabular}{l c c c p{4.5cm}}
\toprule
\textbf{Micro-Sector} & \textbf{Semáforo} & \textbf{Riesgo ML (\%)} & \textbf{Caudal ($Q$)} & \textbf{Diagnóstico Temporal / Terreno} \\
\midrule
\rowcolor{lightbg}
Quebrada El Arrayán Costero & \textbf{\textcolor{alertred}{ALERTA ROJA}} & \textbf{70.3\%} & 4.21 mm & punto máximo transcurrido (10:30 hrs - Calma) \\
Las Rojas \& Entrada Precordillera & \textbf{\textcolor{alertwarning}{ALERTA AMARILLA}} & \textbf{55.0\%} & 2.42 mm & punto máximo transcurrido (10:30 hrs - Drenaje) \\
\rowcolor{lightbg}
Santa Gracia Alta / Pelícano & \textbf{\textcolor{alertwarning}{ALERTA AMARILLA}} & \textbf{52.9\%} & 2.17 mm & punto máximo transcurrido (10:00 hrs - Calma) \\
Las Compañías Alta \& Villa Lambert & \textbf{\textcolor{alertwarning}{ALERTA AMARILLA}} & \textbf{35.2\%} & 1.27 mm & Anegamiento de calles reportado (NLP Activo) \\
\rowcolor{lightbg}
Centro Histórico / Damero & \textbf{\textcolor{alertgreen}{VERDE ESTABLE}} & \textbf{13.5\%} & 0.20 mm & Transitabilidad Normal (Llovizna 1.2 mm) \\
Avenida del Mar / Costero & \textbf{\textcolor{alertgreen}{VERDE ESTABLE}} & \textbf{13.0\%} & 0.18 mm & Borde costero sin anegamiento pluvial \\
\bottomrule
\end{tabular}
\caption{Resultados Obtenidos por el Modelo ML en la Fase de Inferencia en Tiempo Real.}
\end{table}

\section{Beneficios, Puntos a Mejorar y Escalabilidad Institucional}

\subsection{Beneficios para Bomberos y Municipalidad}
\begin{itemize}[leftmargin=*]
    \item \textbf{Zonificación Táctica de Emergencias:} Permite desplegar carros de Bomberos y cuadrillas municipales hacia las zonas con riesgo real (ej: Quebrada El Arrayán o Las Compañías) sin distraer recursos en sectores urbanos estables.
    \item \textbf{Anticipación por Tiempo de Concentración ($T_c$):} Otorga ventanas de prevención diferenciadas de 1.5h en precordillera y de 3.5h a 4h para colectores bajos de la ciudad.
    \item \textbf{Sincronización con SENAPRED y Medios:} Valida de forma automática las Alertas SAE y los reportes ciudadanos en redes y periódicos locales.
\end{itemize}

\subsection{Puntos de Mejora Técnicos}
\begin{enumerate}[leftmargin=*]
    \item \textbf{Sensores de Nivel IoT en Cauces:} Instalar sensores ultrasónicos de bajo costo en las pasarelas del Río Elqui (Islón, Altovalsol, Las Rojas) para telemetría directa de altura de cauce.
    \item \textbf{Refinamiento de Grilla Pluvial Urbana:} Aumentar la resolución a grillas de 100 metros en pasos bajo nivel de Ruta 5 (Km 490 a 500) y arterias principales (Av. Balmaceda, Cuatro Esquinas).
    \item \textbf{API Key Dedicada CEAZAMET:} Formalizar un convenio de acceso web-service directo con el instituto de investigación CEAZA para evitar caídas de scraping durante tormentas.
\end{enumerate}

\subsection{Escalabilidad Institucional}
El Proyecto Júpiter está construido bajo una arquitectura modular en contenedor Docker/Python accesible vía API REST JSON (\texttt{http://localhost:8000/api/scan}):
\begin{itemize}[leftmargin=*]
    \item \textbf{Integración a la Central de Alarmas de Bomberos:} Exportación de capas GIS en formato GeoJSON/KML para ser visualizados directamente en las pantallas del Centro de Despacho de Bomberos de La Serena.
    \item \textbf{Integración a la COGRID Municipal:} Envío automático de resúmenes ejecutivos vía WhatsApp/Telegram a la Dirección de Gestión de Riesgos de Desastres de la Municipalidad.
    \item \textbf{Expansión Regional:} El modelo puede replicarse en menos de 48 horas para las comunas vecinas de Coquimbo, Vicuña, Paihuano y Andacollo.
\end{itemize}

\section{Conclusión}
La Fase de Prueba Real demostró que el Proyecto Júpiter pasó de ser un modelo teórico a una herramienta de inteligencia operativa robusta, capaz de auto-corregirse ante contingencias y entregar información de altísima utilidad táctica para salvar vidas y proteger a la comunidad de La Serena.

\vspace{8mm}
\centering
\textbf{Proyecto Júpiter — Inteligencia Hidrometeorológica para La Serena}\\
\small Región de Coquimbo, Chile — 2026

\end{document}
